\chapter{利未記}
\begin{figure}[H]
\centering
\includegraphics[width=6.in]{ZongJie/figs/AMoXiWuJing/LiWeiJi}
\end{figure}

\section{節期的規定总结}
\subsection{節期的規定总结}
\begin{table}[H] 
\centering
\begin{tabular}{l|l|l|l|l|l}\hline\hline
 節期      &   時間   &  聖會 &    獻的祭物&圣经     \\ \hline
每天& 早晚献祭 &  无&2隻1歲公羊羔+1/10伊法細麵+1 1/4欣油 &民28:1-8\\ \hline
安息日& 每週第七日 &  有&2只没有残疾,一岁的公羊羔+素祭+奠祭 &民28:9-10\\ \hline
月朔& 每月第一日 &-&2只公牛犊,1只公绵羊,7只公羊羔+素祭+奠祭+1只公山羊 &民28:11-15\\ \hline
逾越節& 正月十四日黃昏&-& &出12:1-11\\ \hline
無酵節&正月十五日&  有&公牛犊2+公绵羊1+一岁的公羊羔7+素祭(调油细面)&民28:17-25\\
      &正月十六日&  -&公牛犊2+公绵羊1+一岁的公羊羔7+素祭(调油细面)&民28:17-25\\
      &正月十七日& - &公牛犊2+公绵羊1+一岁的公羊羔7+素祭(调油细面)&民28:17-25 \\
      &正月十八日& -  &公牛犊2+公绵羊1+一岁的公羊羔7+素祭(调油细面)&民28:17-25\\
      &正月十九日& -&公牛犊2+公绵羊1+一岁的公羊羔7+素祭(调油细面)&民28:17-25  \\
      &正月二十日&  -  &公牛犊2+公绵羊1+一岁的公羊羔7+素祭(调油细面)&民28:17-25  \\          
      &正月二十一日&  有&公牛犊2+公绵羊1+一岁的公羊羔7+素祭(调油细面)&民28:17-25\\ \hline
初熟節&莊家初熟& - &初熟莊稼1捆+公綿羊羔1&民28:26-31  \\ %\cline{4-4}
&安息日次日&  &調油細麵2/10伊法+酒1.25欣& \\ \hline
五旬節&初熟節到第七安息& 有 &細麵加酵餅2+羊羔7+公牛犢1&\\% \hline
&日次日(50天)&  &公綿羊2+素祭+奠祭 &\\ \hline
吹角節&七月初一日& 有&  1只公牛犊、1只公绵羊和7只公羊羔+素祭&民29:2-4,8\\ \hline
贖罪日&七月初十日& 有 &1公牛犊+1公绵羊+7只一岁的公羊羔+素祭+奠祭&民29:7-11\\ \hline
住棚節&七月十五日& 有 & 美好樹上的果子、棕樹，柳樹枝+祭物&利23:33-43\\
&&  &13公牛犊、2公绵羊、14只公羊羔+素祭+奠祭+1只公山羊&民29:12-16\\
      &七月十六日& - &12公牛犊、2公绵羊、14只公羊羔+素祭+奠祭+1只公山羊&民29:17-19 \\
      &七月十七日& - &11公牛犊、2公绵羊、14只公羊羔+素祭+奠祭+1只公山羊&民29:20-22 \\
      &七月十八日& - &10公牛犊、2公绵羊、14只公羊羔+素祭+奠祭+1只公山羊&民29:23-25\\
      &七月十九日& - &9公牛犊、2公绵羊、14只公羊羔+素祭+奠祭+1只公山羊&民29:26-28 \\     
      &七月二十日& - &8公牛犊、2公绵羊、14只公羊羔+素祭+奠祭+1只公山羊&民29:29-31 \\
      &七月二十一日& - &7公牛犊、2公绵羊、14只公羊羔+素祭+奠祭+1只公山羊&民29:32-34 \\
      &七月二十二日& 有 &1公牛犊、2公绵羊、14只公羊羔+素祭+奠祭+1只公山羊&民29:35-38  \\  \hline                    
\end{tabular}
\caption{節期的規定}
\end{table}

\begin{table}[H] 
\centering
\begin{tabular}{l|l|l|l|l}
\hline
\hline
 節期      &   時間   &  聖會 &    獻的祭物     \\ \hline
安息日& 每週第七日 &  有 \\ \hline
逾越節& 正月十四日黃昏的時候 &- \\ \hline
無酵節&正月十五日(第一日)&  有&\\
      &正月十六日&  -&\\
      &正月十七日& - & \\
      &正月十八日& -  &\\
      &正月十九日& -&  \\
      &正月二十日&  -  &  \\          
      &正月二十一日(第七日)&  有&\\ \hline
初熟節&莊家初熟安息日的次日& - &初熟莊稼1捆+公綿羊羔1+調油細麵2/10伊法+酒1.25欣 \\ \hline
五旬節&初熟節到第七安息日次日(50天)& 有 &細麵加酵餅2+羊羔7+公牛犢1+公綿羊2+素祭+奠祭 \\ \hline
吹角節&七月初一日& 有 & 要將火祭獻給耶和華\\ \hline
贖罪日&七月初十日& 有 &要將火祭獻給耶和華\\ \hline
住棚節&七月十五日(第一日)& 有 & 美好樹上的果子、棕樹，柳樹枝+祭物\\
      &七月十六日& - & \\
      &七月十七日& - & \\
      &七月十八日& - & \\
      &七月十九日& - &  \\     
      &七月二十日& - & \\
      &七月二十一日& - &  \\
      &七月二十二日(第八日)& 有 &   \\  \hline                    
\end{tabular}
\caption{節期的規定}
\end{table}





\subsection{承接聖職之禮的总结 8}
  \begin{tikzpicture}[%
    >=stealth,
    node distance=3cm,
    on grid,
    auto
  ]
    \node[state] (A)              {A};
    \node        (B) [right of=A,fill=blue!25,text width=2cm]{水洗亚伦和他兒子们 };;
    \path[->] (A) edge (B);
  \end{tikzpicture}


\sffamily\begin{tikzpicture}
  % Place all element in a matrix of nodes, called m
  % By default all nodes are rectangles with round corners
  % but some special sytles are defined also
  \matrix (m) [matrix of nodes, 
    column sep=5mm,
    row sep=1cm,
    nodes={draw, % General options for all nodes
      line width=1pt,
      anchor=center, 
      text centered,
      rounded corners,
      minimum width=1.5cm, minimum height=8mm
    }, 
    % Define styles for some special nodes
    right iso/.style={isosceles triangle,scale=0.5,sharp corners, anchor=center, xshift=-4mm},
    left iso/.style={right iso, rotate=180, xshift=-8mm},
    txt/.style={text width=1.5cm,anchor=center},
    ellip/.style={ellipse,scale=0.5},
    empty/.style={draw=none}
    ]   
  {
  % Second row of symbols
  % m-1-1
    水洗亚伦和他兒子们 
  & % m-1-2
    給亞倫穿聖衣
  & % m-1-3
    膏抹帳幕和其中所有  
   & % m-1-4
    比哈希錄  
   & % m-1-5
    |[draw=orange!80!white, ultra thick]| \textbf{密奪}   
   & % m-1-6
    |[coordinate, xshift=-1cm]|  
  \\  
  % Third row of symbols
    % m-2-1  
   瑪撒（探米）或叫米利巴（爭鬧）
  & % m-2-2試探
    |[draw=orange!80!white, ultra thick]| \textbf{汛的曠野} 
  & % m-2-3    
    以琳 
  & % m-2-4  
        |[draw=orange!80!white, ultra thick]| \textbf{瑪拉}       
  & % m-2-5  
    書珥曠野  
  & % m-2-6  (no symbol here, only a point to draw the path)
    |[coordinate, xshift=-1cm]| 
  \\  
利非訂
  & % m-3-1試探
  |[draw=orange!80!white, ultra thick]| \textbf{西乃曠野} 
  & % m-3-1試探
     |[coordinate, xshift=-1cm]| 
      \\     
  };  % End of matrix

  % Now, connect all nodes in a chain.
  % The names of the nodes are automatically generated in the previous matrix. Since the
  % matrix was named ``m'', all nodes have the name m-row-column
  { [start chain,every on chain/.style={join}, every join/.style={line width=1pt}]
    %\chainin (m-1-1);
  	 \chainin (m-1-1);
    \chainin (m-1-2);
    \chainin (m-1-3); 
    \chainin (m-1-4); 
    \chainin (m-1-5); 
    \path (m-1-5) edge node [below] {\small{3天}} (m-2-5) ; 
    \chainin (m-2-5);     
    \chainin (m-2-4);
    \chainin (m-2-3); 
   \chainin (m-2-2);  
    \chainin (m-2-1);  
   \chainin (m-3-1); 
       \path (m-3-1) edge node [below] {\small{葉忒羅}} (m-3-2) ;  
     \chainin (m-3-2);             
    };

                rotate=-90, minimum height=7mm, minimum width=1.3cm, 
                single arrow head extend=1.2mm, single arrow tip angle=120] {};
 
  % Define the style for the blue dotted boxes
  \tikzset{white dotted/.style={draw=white!50!white, line width=1pt,
                               dash pattern=on 1pt off 4pt on 6pt off 4pt,
                                inner sep=4mm, rectangle, rounded corners}};

  % Finally the blue dotted boxes are drawn as nodes fitted to other nodes
 % \node (first dotted box) [blue dotted, 
  %                          fit = (MOD text) (m-2-1) (m-2-4)] {};
  \node (second dotted box) [white dotted,
                            fit = (m-3-2)] {};
  \node at (second dotted box.south) [above, inner sep=1.mm] { \small{頒布律法，建會幕}};
  \node (second dotted box) [white dotted,
                            fit = (m-2-1)] {};
  \node at (second dotted box.south) [above, inner sep=1.mm] { \small{摩西擊打磐石}};

   \node (second dotted box) [white dotted,
                            fit = (m-3-1)] {};
  \node at (second dotted box.south) [above, inner sep=1.mm] { \small{與亞瑪力爭戰}};
 
  \node (second dotted box) [white dotted,
                            fit = (m-2-4)] {};
  \node at (second dotted box.south) [above, inner sep=1.mm] { \small{苦水變甜}};

  \node (second dotted box) [white dotted,
                            fit = (m-2-2)] {};
  \node at (second dotted box.south) [above, inner sep=1.mm] { \small{降鵪鶉和嗎哪}};


  % Define the style for the blue dotted boxes
  \tikzset{blue dotted/.style={draw=blue!50!white, line width=1pt,
                               dash pattern=on 1pt off 4pt on 6pt off 4pt,
                                inner sep=4mm, rectangle, rounded corners}};

  % Finally the blue dotted boxes are drawn as nodes fitted to other nodes
 % \node (first dotted box) [blue dotted, 
  %                          fit = (MOD text) (m-2-1) (m-2-4)] {};
  \node (second dotted box) [blue dotted,
                            fit = (m-1-1) (m-1-2) (m-1-3) (m-1-4) (m-1-5)] {};

  % Since these boxes are nodes, it is easy to put text above or below them                          
  %\node at (first dotted box.north) [above, inner sep=3mm] {\textbf{Transmitter}};
  %\node at (second dotted box.south) [below, inner sep=3mm] { \tiny{過紅海前}};
  \node at (second dotted box.south) [above, inner sep=1.mm] { \small{過紅海前}};

\end{tikzpicture}  

  
%\noindent
摩西用水洗了亞倫和他兒子們\\
給亞倫穿祭司服裝\\
摩西用膏油抹帳幕和其中所有的，又抹了壇和壇的一切器皿，並洗濯盆和盆座\\
把膏油倒在亞倫的頭上膏他\\
給亞倫的兒子們穿上內袍，束上腰帶，包上裹頭巾\\
摩西為祭司獻贖罪祭的公牛 \\
摩西奉上燔祭的公綿羊，是獻給耶和華的火祭\\
他又奉上第二隻公綿羊，就是承接聖職之禮的羊。又取一個無酵餅，一個油餅，一個薄餅，都放在脂油和右腿上。亞倫和他兒子用手作搖祭，在耶和華面前搖一搖。摩西從他們的手上拿下來，燒在壇上的燔祭上，都是為承接聖職獻給耶和華馨香的火祭。\\
摩西拿羊的胸作為搖祭，在耶和華面前搖一搖，是承接聖職之禮，歸摩西的分。\\
摩西取點膏油和壇上的血，彈在亞倫和他的衣服上，並他兒子和他兒子的衣服上，使他和他們的衣服一同成聖。\\
亞倫和他兒子把肉煮在會幕門口，又吃承接聖職筐子裡的餅\\
他們七天不可出會幕的門，晝夜住在會幕門口等到承接聖職的日子滿了，為他們贖罪。



\subsection{亞倫獻祭总结 9}
%\noindent
1）亞倫為自己獻祭\\
獻贖罪祭-公牛犢（1）\\
獻燔祭-公綿羊（1）\\
%\noindent
2）亞倫為百姓獻贖祭\\
獻贖罪祭的公山羊（1）\\
獻燔祭一歲沒有殘疾的牛犢（1）和綿羊羔（1）\\
平安祭公牛（1）和公綿羊（1）。胸和右腿，亞倫當作搖祭\\
奉上調油的素祭，從其中取一滿把，燒在壇上


\subsection{可吃,不可吃的食物总结 11}
%%\noindent
\begin{table}[h]%{r}{4.2cm}
\vspace*{-\baselineskip}
\begin{tabular}{ll l}
\hline
\hline
种类 &可以吃& 不可吃的 \\\hline
地上一切走獸 &凡蹄分兩瓣、倒嚼的走獸 & 倒嚼不分蹄：駱駝，沙番，兔子。分蹄不倒嚼；豬 \\\hline
水裡海裡河裡&有翅有鱗的&無翅無鱗的，你們都當以為可憎 \\\hline
雀鳥中& &吃腐肉类——鵰，狗頭鵰，紅頭鵰，鷂鷹，小鷹與其類,烏鴉與其類   \\
& &夜行类——夜鷹，魚鷹，鷹與其類，鸮鳥，貓頭鷹，角鴟，  \\
& &地上不飞的鸟——鴕鳥，鸕鶿，鵜鶘（吃多了就不飞）；  两栖类鸟——蝙蝠 \\
& &水上不常飞的鸟——禿鵰，鸛，鷺鷥與其類，戴鵀（在粪上觅食）  \\\hline
昆虫类 &有翅膀有足在地上蹦跳的&凡有翅膀用四足爬行的物，你們都當以為可憎   \\\hline
 地上爬物&&鼬鼠，鼫鼠，蜥蜴與其類，壁虎，龍子，守宮，蛇醫，蝘蜓。 \\\hline
\end{tabular}
\caption{可吃,不可吃的食物总结}
\end{table}



\subsection{各項器皿，用具，食物，水的潔淨規定总结 11}
%\noindent
木器、衣服、皮子、口袋，要放在水中，必不潔淨到晚上，到晚上才潔淨了。\\
瓦器，爐子，鍋臺沾了不潔淨要被打破。\\
泉源或是聚水的池子仍是潔淨\\
若是死的，有一點掉在要種的子粒上，子粒仍是潔淨；若水已經澆在子粒上，那死的有一點掉在上頭，這子粒就與你們不潔淨

\subsubsection{大麻風病的鑑定总结 13}
%\noindent
 1）肉皮癬，或火斑與大麻風的辨別\\
 2）皮上白癤與大麻風的辨別\\
 3）皮肉上長瘡之處又起了白癤、或是白中帶紅的火斑,與大麻風的辨別\\
 4）皮肉上火毒成了火斑與大麻風的辨別\\
 5）在頭上（頭疥）災病，或是男人鬍鬚上大麻風的辨別\\
 6）皮肉上白火斑，白癬  與大麻風的辨別\\
 7）頭禿，與大麻風的辨別





















